\documentclass[10pt,twocolumn,letterpaper]{article}
\usepackage[T1]{fontenc}
\usepackage[review]{cvpr}
\usepackage{amsmath,amssymb,booktabs,graphicx,placeins}
\usepackage[breaklinks,colorlinks,allcolors=blue]{hyperref}
\def\confName{CVPR}
\def\confYear{2027}
\def\paperID{TBD}
\title{What Inertial Recordings Reveal About Furniture Layout}
\author{Anonymous submission}
\setlength{\dbltextfloatsep}{10pt plus 3pt minus 3pt}
\setlength{\textfloatsep}{9pt plus 3pt minus 3pt}
\setlength{\dblfloatsep}{8pt plus 2pt minus 2pt}
\setlength{\floatsep}{7pt plus 2pt minus 2pt}
\setcounter{topnumber}{3}
\setcounter{dbltopnumber}{2}
\setcounter{totalnumber}{4}
\renewcommand{\topfraction}{.92}
\renewcommand{\dbltopfraction}{.92}
\renewcommand{\textfraction}{.08}
\renewcommand{\floatpagefraction}{.72}
\renewcommand{\dblfloatpagefraction}{.72}
\begin{document}
\maketitle
\begin{abstract}
We quantify how wearable activity discloses furniture positions under
known room outlines, starts, inventories and nominal dimensions.
Separate evaluations measure ideal geometric evidence, acquisition from
synthetic phone and watch IMUs, and real-room feasibility. Across 10,000
controlled layouts, dispersing contacts at the same walking route and
fourteen interactions raises whole-room centre recovery from 39.78\% to
64.70\%, a paired gain of 24.92 percentage points with a room-bootstrap
95\% interval of {[}24.00, 25.84{]} and a mean room-level
instance-position F1 of 90.43\%. The same intervention end-to-end through
noisy phone and wrist streams, trajectory refiners and a frozen decoder
reproduces on 64 rooms that no earlier arm uses: 48/64 against 30/64, a
gain of 28.12 points ({[}15.62, 40.62{]}, exact McNemar p = 1.2e-4).
Declared side information alone recovers 0/32 records, and an
independent sampler over 64,000 declared-only layouts over 32 rooms
reaches F1 0.088 with four exact rooms; selecting the best draw with
reference layouts raises F1 to 0.737 with 4/32 exact, while the inertial
chain reaches 0.801 with 21/32. Fixed-record subset controls show that
selection and refinement, not contact count alone, set the operating
point. A production replay localises four objects in one real room with
centre errors of 0.09--0.23 m. Contact placement is therefore a
first-order lever that survives the device chain, and declared priors
alone do not substitute for it.
\end{abstract}
\section{Introduction}\label{introduction}

Wearable inertial recordings reflect movement and its spatial
constraints. A walk passes through free space, sitting locates a seat
region, and opening a cabinet constrains an operation surface. Given a
room's dimensions and inventory, these observations can narrow furniture
positions. Quantifying that disclosure requires specifying the
observations and available side information.

Existing systems recover layouts from trajectories or synthesise scenes
compatible with motion. Walk2Map learns trajectory-to-plan relationships
\cite{r1}, CoGuide studies spatial inverse problems and varies
trajectory density and noise \cite{r2}, and SUMMON synthesises
contact-compatible objects \cite{r3}. Wearable systems extend motion
understanding to scenes and interactions \cite{r4,r12}. We examine a
more specific experimental question. At the same walking route and
number of interactions, how much does changing the spatial placement of
contacts affect furniture-position recovery?

The answer depends on the measurement protocol. In our controlled
geometric experiment, moving fourteen contacts from repeated to
dispersed locations changes whole-room recovery by 24.92 percentage
points. We then rebuilt that intervention inside the device chain, with
identical rooms, route, fourteen contacts, sensor seeds, trajectory
refiners and decoder, and the placement protocol as the only difference.
On 64 rooms that no earlier arm uses, the effect persists at 48/64
against 30/64, a gain of 28.12 points, so the geometric lever is not an
artefact of ideal trajectories. A twelve-room development block, reported
separately, is too small to resolve a difference of that sign. Declared
near-wall structure and finite decoding also shift recovery at a given
activity budget, but declared priors alone do not reproduce it: across
64,000 sampled layouts the best attainable F1 is 0.088. These
measurements define conditional operating points for a scene family,
side-information package, interaction policy and decoder. Walking
coverage uses reference free space as its denominator, so these maps
describe controlled activity budgets and require additional evidence
before they can serve as a runtime stopping rule.

The device experiment then asks which of these spatial observations a
frozen inference chain can extract. On two synthetic batches, removing
inertial evidence while retaining the same finite decoder and declared
side information gives 0/32 whole-room successes, compared with 17/32
and 21/32 for the corresponding raw-input arms. This evidence ablation
evaluates the contribution of observations to that estimator. Its
interpretation is bounded by the estimator's prior-only search and
ordering. Within fixed recordings, further controls vary which contacts
reach the refiners and decoder, making evidence selection a measurable
part of the pipeline.

We contribute a paired contact-placement intervention over 10,000
controlled layouts, an end-to-end replication of that intervention on 64
unseen rooms of the device chain, conditional budget maps retaining
partial and whole-room endpoints, and device controls separating
acquisition, feasibility and selection. The three layers retain separate
chains and denominators. Device training and testing use the same
generation family, and real-input feasibility uses one development room.

\begin{figure*}[t]
\centering
\includegraphics[width=\textwidth]{../../figures/cvpr/fig_framework031.pdf}
\caption{The framework links activity, acquired evidence and position disclosure,
with higher F1 and whole-room counts indicating better recovery. Device
scores refer to 32 synthetic records, while room drawings and waveforms
are illustrative. The three datasets retain separate chains and
denominators, the watch supplies contact and timing evidence, and
physical clock synchronisation is unverified.}
\end{figure*}

\section{Related work}\label{related-work}

\textbf{Scene recovery from movement.} Walk2Map learns floor plans and
coarse furniture regions from walk trajectories \cite{r1}. CoGuide
formulates spatial inference using contrastive diffusion guidance and
studies sparse and noisy trajectories \cite{r2}. SUMMON uses human
motion and contact to synthesise plausible scene objects \cite{r3}.
These studies motivate our use of movement as spatial evidence. Our
paired intervention holds route, interaction count and decoder fixed
while changing contact placement.

\textbf{Wearable interaction and localisation.} IMU4D estimates body
motion and coarse scene structure from wearable IMUs \cite{r4}. IMU-HOI
uses body- and object-mounted IMUs for coherent human--object motion
\cite{r12}. RoNIN addresses inertial odometry \cite{r6}, while
ActionSLAM uses actions as localisation landmarks \cite{r5}. Collective
landmark mapping also addresses ambiguities between class labels and
physical instances \cite{r13}. We retain unknown event-to-instance
associations, use phone and watch inertial channels under a declared
carrying protocol, and place no sensor on furniture.

\textbf{Position disclosure as a measurement task.} We evaluate the
furniture positions recoverable by a specified estimator under stated
observations and side information. Partial instance recovery, complete
room recovery and footprint agreement answer different questions. This
distinction matters when a plausible output contains a misplaced object,
or when a finite decoder retains useful individual candidates but
selects an incorrect joint layout.

\section{Task, evidence and
evaluation}\label{task-evidence-and-evaluation}

\subsection{Conditional activity
budgets}\label{conditional-activity-budgets}

Each controlled rectangular room contains one bed and one or two
instances of each desk, storage and chair category, giving four to seven
target instances. The 10,000-room set contains 55,000 target items. The
analyst receives the room outline, starting position, furniture
categories, counts and nominal dimensions. Furniture positions, absolute
heading and event-to-instance associations are inferred. Axis alignment,
interaction semantics and any near-wall assumptions are stated with the
corresponding experiment.

Walking coverage measures the fraction of reference traversable grid
points within 0.25 m of the walked path. The route construction uses an
approximately 0.20 m grid and 0.15 m body clearance. Let G\_free be the
reference free-space grid and T the sampled trajectory. We report

\[c(T)=\frac{|\{g\in G_{\mathrm{free}}:\min_{p\in T}\|g-p\|_2\leq0.25\,\mathrm{m}\}|}{|G_{\mathrm{free}}|}.\]

Coverage is an experimental design and evaluation variable. Reference
furniture geometry plans and scores the motion, while the inference
interface receives the declared side information and observations. The
total interaction budget n follows a chair--storage--bed cycle, rotating
among same-category instances. Repeated and dispersed arms differ in
contact location. Desks have no direct contact channel and are
constrained through walking, side information and relations to seated
events.

\subsection{Endpoints}\label{endpoints}

Predictions and references are matched one to one within category, with
a Euclidean centre tolerance of 0.5 m. For room r, we compute

\[\begin{gathered}F_{1,r}=\frac{2TP_r}{2TP_r+FP_r+FN_r},\\[4pt]H_r=\mathbf{1}[TP_r=|\widehat{O}_r|=|O_r|].\end{gathered}\]

We average each endpoint over rooms. Whole-room centre success requires
all instances and no extras. An empty output receives zero on both
endpoints. Decoder solve status and the separate endpoint requiring
every matched footprint to reach IoU of at least 0.5 are also retained.

For a declared protocol Pi, the budget map reports the empirical mean F1
and mean H at each measured (c,n) cell. Pi includes the room family,
side information, interaction policy, acquisition settings and decoder.
This notation describes the measured operating points. The finite
candidate domains and search limits remain part of their interpretation.

\subsection{Input chains and inference
units}\label{input-chains-and-inference-units}

Layer 1 receives ideal trajectories and class-level contact coordinates.
Layer 2 starts from synthetic phone acceleration and angular velocity at
120 Hz and watch acceleration at approximately 16.7 Hz. Its
segment-displacement regressor is ExtraTrees \cite{r10} with 300 trees,
minimum leaf size two and 148 features, trained on 7,723 segments from
historical rooms 0--19. Model selection uses records 11000--11007, while
primary confirmation records 11100--11131 are excluded from that
selection. Training and evaluation share the generation family. Some
evaluation geometries appeared in earlier project stages, so the holdout
concerns the stated model-selection process and newly generated long
records.

The phone activity model uses a frozen random forest \cite{r9}, and a
watch-only cabinet channel supplies wrist events. Under the declared
pose-to-body protocol, sitting anchors seat regions, lying supports the
bed footprint, and cabinet events refer to an operation face with
nominal reach 0.60 m. Unknown point and plane revisits regularise
trajectory drift. Directional finite decoding uses cap 128, contact
padding 0.30 m and walking tolerance 0.20 m. Later diagnostic chains
additionally apply storage-wall trajectory repair and gated gyro-bias
tracking. Each arm identifies its own chain. The real-room layer retains
the original RoNIN production entry point and its room-specific
relations.

Rooms are the inferential units. Binary paired comparisons use exact
two-sided McNemar tests. Proportions use pointwise Wilson 95\% intervals
\cite{r11}, and paired effects use room-bootstrap intervals where
stated. In the replicated grid, each room contributes zero to three
successes from three handset-noise realisations, with wrist streams
fixed. Exploratory Wilcoxon tests are two-sided and discard zero
differences. Six simple effects are reported with Holm sensitivity
results in the supplement. A direct post-hoc room-level interaction
contrast uses sign flips, assuming independent rooms and symmetric
contrasts under the null.

\section{Geometric evidence and conditional
budgets}\label{geometric-evidence-and-conditional-budgets}

\subsection{Contact placement at fixed route and
count}\label{contact-placement-at-fixed-route-and-count}

Figure 2 retains all 43 measured cells. Mean room-level position F1
ranges from 4.56\% to 90.43\%, while whole-room recovery over those
cells ranges from 0\% to 64.70\%. The heatmap and three-dimensional
lattice encode the same F1 values. Five unmeasured cells are marked
explicitly.

At the same route and fourteen interactions, repeated contact placement
recovers 39.78\% of rooms and dispersed placement 64.70\%. The paired
difference is 24.92 percentage points, with room-bootstrap 95\% interval
{[}24.00, 25.84{]}. Mean distinct contact coordinates increase from 4.0
to 10.5. This intervention isolates contact placement within the
controlled geometric protocol. At full coverage and fourteen dispersed
interactions, instance-position F1 is 90.43\%, but only 740/10,000 rooms
meet the separate whole-room footprint-IoU endpoint. The distinction
limits conclusions about shapes and orientations.

\begin{figure}[t]
\centering
\includegraphics[width=\columnwidth]{../../figures/cvpr/fig_budget032.pdf}
\caption{Activity budgets set the operating level under a fixed geometric
protocol. Panel a gives room-macro position F1 for all 43 measured cells,
each resting on 10,000 controlled rooms, with hatching for the five
unmeasured cells. Panel b replots the full-coverage column against
interaction count, where three more interactions buy 3.3 points once
coverage is complete. Values are F1, not whole-room counts.}
\end{figure}

\subsection{Dependence on priors, decoding and scene
family}\label{dependence-on-priors-decoding-and-scene-family}

A wall-property audit of 1,000 historical generated rooms finds every
desk within 0.15 m of a wall, with median gap 0.104 m. Adding this
declared property to the geometric candidate pool raises whole-room
recovery from 651 to 826 on the same 1,000 development rooms and desk
recall from 75.87\% to 87.27\%. A new same-family 300-room set gives
66.0\% and 84.0\%. On a separate 1,000-room confirmation set,
post64\_wall40k recovers 851 rooms and depth\_portfolio 945 at the same
maximum total budget of 40,000 nodes. The paired gain is 9.4 percentage
points with bootstrap 95\% interval {[}7.6, 11.3{]}.

With the declared desk-wall prior in the development configuration,
measured operating points are 64.7\% at 95\% coverage/14 interactions,
70.0\% at 100\%/8, 82.6\% at 100\%/14 and 85.1\% at 100\%/20. The 82.6\%
point is conditional on that configuration and room family. The
corresponding outline-decoder result over 10,000 rooms is 64.70\%,
increasing to at most 66.63\% at the separately tested 23-interaction
budget. All sets and configurations are kept distinct.

Scene transfer further changes the operating point. Of 25 exported iTHOR
scenes \cite{r7}, fifteen contain the target categories and thirteen
pass the geometry screen. The earlier 11/11 result uses an additionally
restricted subset. On 803 separately included rectangular ProcTHOR rooms
\cite{r8}, 95\% coverage and five interactions yield 350 centre
successes, or 43.6\%. That earlier protocol selects one representative
per target role and maps Dresser to storage, with idealised
storage-surface contact semantics. Its endpoint differs from the
variable-inventory experiment.

\section{What raw device streams
contribute}\label{what-raw-device-streams-contribute}

\subsection{Observation and pose
interventions}\label{observation-and-pose-interventions}

The primary raw joint arm recovers 17/32 rooms (53.13\%, Wilson 95\%
{[}36.45, 69.13{]}) with mean instance-position F1 0.7275. The older
front end recovers 11/32. Their paired outcomes are nine both
successful, eight new-only successes, two old-only successes and
thirteen both unsuccessful. The 18.75-point difference has paired
bootstrap 95\% interval {[}0, 37.5{]} and exact McNemar p = 0.1094.

Removing wrist evidence gives 9/32, with eight joint-only successes and
no wrist-absent-only success (unadjusted diagnostic p = 0.0078).
Shifting the wrist recording by ten seconds gives 0/32. Replacing
estimated pose with reference pose, while retaining predicted actions,
gives 26/32 and F1 0.9536. This diagnostic changes the pose-dependent
refinements and finite decoding together. Six rooms still fail to meet
the centre endpoint.

\textbf{Evidence ablation with the same decoder.} Feeding only outline,
start and inventory to the directional decoder, with cap 128 and support
weight 0.3, gives 0/32 whole-room successes and mean position F1 0.0915.
All 32 runs return layouts. Against the corresponding raw arm, the
paired difference is 17 rooms, with bootstrap 95\% interval {[}34.4,
68.8{]} percentage points and exact McNemar p = 1.5e-5. A second 32-room
set gives 0/32 and F1 0.0793, compared with its stored raw arm's 21/32
and 0.8013 (p = 9.5e-7). This post-hoc ablation retains one finite
estimator's candidate ordering and search behaviour.

\textbf{Limit of the estimator or of the priors?} An independent sampler
draws 2,000 layouts per room from the declared inventory, dimensions,
non-overlap, wall-clearance and support constraints, and reproduces the
ablation: F1 0.0878 with 4/64,000 exact recoveries over the same 32
rooms. The failure is therefore not specific to one decoder. Selecting
the best of the 2,000 draws per room with reference layouts gives F1
0.7371 and 4/32 exact rooms, and that ceiling falls as inventories grow:
0.950 for four items (4/5 rooms), 0.725 for five (0/13), 0.716 for six
(0/14) and 0.595 for seven (0/12). Above five items, 78,000 draws contain
no exact room (Fig. 5). Declared priors in this family support
approximate placement at best, and the inertial chain's 21/32 on the same
rooms is what supplies exactness. The ceiling belongs to one sampling
family rather than to all prior-only inference.

\begin{figure*}[t]
\centering
\includegraphics[width=\textwidth]{../../figures/cvpr/fig_prior033.pdf}
\caption{What declared side information can recover without inertial
evidence. Panel a compares one finite decoder fed only outline, start and
inventory with an independent 2,000-draw sampler on the same 32 rooms,
together with a sampler draw selected using reference layouts and the
inertial chain. Panel b reports that truth-assisted ceiling against
inventory size, where it drops from 0.950 at four items to 0.595 at
seven. Both panels share one room set and one declared input package.}
\end{figure*}

\begin{table}[t]
\centering\small
\caption{Primary 32-record controls. Higher position F1 and whole-room counts indicate better recovery. All rows retain the same 0.5 m centre endpoint.}
\begin{tabular}{@{}p{.50\linewidth}rr@{}}\toprule
Arm & Centres & F1 \\\midrule
Declared side information only & 0/32 & 0.0915 \\
Older trajectory front end & 11/32 & 0.4375 \\
Raw phone and watch & 17/32 & 0.7275 \\
Reference pose, predicted actions & 26/32 & 0.9536 \\
\bottomrule\end{tabular}
\end{table}

\subsection{Placement transfers to the device
chain}\label{placement-transfers-to-the-device-chain}

The geometric intervention changes only where contacts land. We rebuilt
it inside the Layer-2 chain: identical room, route, fourteen contacts,
handset and wrist seeds, refiners and decoder, with the placement
protocol as the only difference. The twelve development rooms give 7/12
(mean F1 0.8135) for dispersed placement and 8/12 (0.8379) for repeated
placement, a difference of $-8.33$ points with exact McNemar p = 1.0.
At that size the block cannot resolve an effect of the geometric
magnitude. A fresh block of 64 rooms, which no earlier arm in this paper
reports on, separates the two protocols: dispersed placement recovers
48/64 (75.0\%, F1 0.9304) and repeated placement 30/64 (46.9\%,
0.8354). Twenty rooms succeed under dispersed placement only, against
two under repeated placement only, giving a paired gain of 28.12 points
with room-bootstrap 95\% interval {[}15.62, 40.62{]} and exact McNemar p
= 1.2e-4 (Fig. 4). Pooling all 76 rooms gives 55/76 against 38/76,
+22.37 points ({[}10.53, 34.21{]}, p = 9.1e-4). The sign and magnitude
match the geometric experiment, and the twelve-room block was
underpowered rather than contradictory. Both device blocks share one
finite chain, so this estimate is conditional on its acquisition and
refinement and does not establish transfer to another decoder.

\begin{figure*}[t]
\centering
\includegraphics[width=\textwidth]{../../figures/cvpr/fig_placement033.pdf}
\caption{Contact placement transfers from the geometric protocol to the
device chain. Panel a gives paired room outcomes for 64 rooms under
dispersed and repeated placement with the same route, fourteen contacts,
seeds, refiners and decoder. Panel b compares the paired effect with the
10,000-room geometric experiment and with the twelve-room development
block, with room-bootstrap 95\% intervals. Dispersed placement gains 28.1
points on the fresh block and loses 8.3 points on the underpowered
twelve-room block.}
\end{figure*}

\subsection{Evidence selection within fixed
recordings}\label{evidence-selection-within-fixed-recordings}

The replicated visit-by-error grid contains twelve rooms, three visit
counts, three handset-error levels and three handset-noise realisations
per condition. Whole-room counts are 18/16/14 at six visits, 21/20/12 at
ten visits and 23/21/11 at fourteen visits, each out of 36 observations
(Fig. 5a). The independent room count remains twelve. Increasing visits
from six to fourteen changes success by +13.89 percentage points at
severity zero and $-8.33$ points at severity one. The direct
difference-in-differences is $-22.22$ points, with exploratory
sign-flip p = 0.125. Visits also lengthen the recording, so interpreting
a separate contact-count effect requires further controls.

\textbf{Holding the recording fixed.} On the same fourteen-visit record
per room, all fourteen qualified contacts recover 7/12 rooms with F1
0.8135. The first six recover 4/12 with 0.7274, dispersed six recover
5/12 with 0.7028, and three random six-contact subsets per room recover
8/36 with 0.6435 (Fig. 5b). A truth-assisted six-contact heuristic gives
4/12 with 0.7960. It accesses reference object centres and is evaluated
as a diagnostic. The same raw recording, front end, walking runs and
decoder settings are retained, while the selected events feed the
contact and wall refiners before decoding. Thus the contrast includes
selection-induced changes in the refined trajectory and heading
hypotheses.

\textbf{A second recording length.} Within the twenty-visit records, all
twenty contacts and the first fourteen each recover 8/12 rooms, with F1
0.7817 and 0.7917. Dispersed and truth-assisted fourteen-contact subsets
also give 8/12, with F1 0.8393 and 0.8671. Equal aggregate hit counts
can involve different successful rooms. The longer records last
1.18--1.26 times as long as the fourteen-visit records. Cross-length
differences therefore combine changed motion, added signal and
refinement effects. The within-record observations establish equality at
the tested budgets, while the separate grid's visit-and-duration
dependence remains unresolved.

\begin{figure*}[t]
\centering
\includegraphics[width=\textwidth]{../../figures/cvpr/fig_evidence032.pdf}
\caption{Contact availability and contact selection under different controls.
Panel a reports whole-room success over 36 observations per cell from
twelve rooms and three handset-noise realisations; the interaction is
exploratory, with p = 0.125. Panel b holds one record per room and varies
only which six contacts reach the refiners and decoder, against all
fourteen and against a truth-assisted diagnostic. Cells are success
counts, not F1.}
\end{figure*}

Ranking contacts by event duration provides a useful negative control.
Keeping the longest eleven, nine and six contacts gives 4/12, 0/12 and
0/12 whole-room successes. Corresponding three-seed random subsets give
19/36, 17/36 and 14/36. Duration ranking retains cabinet events and
removes all lying contacts at six and nine contacts. The synthetic event
schedule makes duration category-dependent, so this result bounds that
score in this chain. Detector confidence remains a separate question.
Complete subset results and class counts are in the supplement.

\textbf{Pinning the class mix.} A second selection experiment removes the
class composition as a free variable. Each room receives six contacts
allocated across seated, lying and cabinet events in proportion to the
available counts, with at least one per present class; all twenty-four
rooms here admit the quota 2/2/2. Under that fixed quota a spatially
dispersed subset beats three same-quota random subsets on a confirmation
block of twelve rooms that the rule never saw, 11700--11711: nine rooms
better, one worse, two equal, mean F1 +0.0899, whole-room hits 7/12
against 16/36. The twelve development rooms agree in sign (+0.0503,
five better, three worse, four equal). The quota itself is not the
source: same-quota random against unrestricted random gives $-0.0389$ on
the confirmation block, although the two are +0.1204 apart on
development. A quota-respecting truth-assisted subset stays below random
in both blocks (0.6806 and 0.7252), so taking the nearest contact per
object is not an upper bound here. Six contacts, however chosen, remain
below the fourteen-contact arm (0.8829 against 0.9603 on the
confirmation block). With twelve rooms per block and one severity, these
sign counts are descriptive (Table S13).

\subsection{Temporal error and remaining inference
failures}\label{temporal-error-and-remaining-inference-failures}

On twelve rooms, jointly scaling handset white noise, bias walk,
scale/misalignment and timestamp jitter gives 7/3/1/0 whole-room
successes at severities 0/0.5/1/2. The severity-0.5 acceleration RMS
difference is 0.0132 m/s\^{}2, but multiple error terms change together.
Per-term controls give 3/12 for bias walk, 5/12 for scale/misalignment
and 7/12 for either white noise or jitter. Per-term RMS values remain
unreported.

At fixed nominal per-axis added RMS of 0.020 m/s\^{}2 for acceleration
and 0.003 rad/s for gyro, temporal structure also matters. Clean,
white-gyro, constant-gyro and slowly varying-gyro conditions recover
10/12, 11/12, 8/12 and 0/12, respectively. Slowly varying acceleration
gives 8/12. Slowly varying gyro errors can prevent the frozen front end
from detecting its required initial stationary block. Supplementary
Figure S1 and Table S6 retain the channel and correction controls.

Gated gyro-bias tracking leaves the unperturbed confirmation arm at
24/32. At severity 0.5 it changes 12/32 to 18/32 (p = 0.03125, Holm p =
0.09375), and at severity one 3/32 to 14/32 (p = 0.0034, Holm p =
0.0137). A second batch changes 13/32 to 16/32 at severity 0.5. Across
the two severity-0.5 batches, 25/64 becomes 34/64 through eleven gains
and two losses (p = 0.0225). For slowly varying gyro error in the
twelve-room diagnostic, gating restores 5/12 while retaining the clean
10/12. These results concern the evaluated error generator and frozen
front end.

Across all 324 replicated rows, 266 return a solution and 156 recover
every centre. The remaining 110 solved rows contain at least one
position error. The 58 unsolved rows comprise 45 proposal conflicts,
eleven exhausted search budgets and two empty proposal pools. These are
repeated observations from twelve rooms. A separate 32-room finite-pool
diagnostic retains a near-truth candidate for every object, yet
restricting candidates using evaluation-time truth yields only 28/32
jointly feasible correct layouts, compared with 24/32 for the deployed
arm. Candidate cap 64 gives 25/23/18 against 24/23/17 across three
batches. Other tested selection criteria lack consistent cross-batch
gains. These diagnostics distinguish candidate retention, joint
compatibility and final selection.

The inventory assumption also affects feasibility. These twelve rooms
contain four to seven instances, averaging 5.67. The correct list
recovers 7/12 with F1 0.8135. Removing one declared storage instance
yields nine proposal conflicts and mean F1 0.1965, while removing one
chair yields five conflicts and F1 0.4349. Mean true positives are 4.58,
1.08 and 2.33, respectively, distinct from mean output counts of 5.25,
1.25 and 2.92. Whole-room success is zero by definition after a
reference instance is removed from the declaration. Feasibility and
matched-instance counts therefore carry the substantive comparison.

\begin{figure*}[t]
\centering
\includegraphics[width=\textwidth]{../../figures/cvpr/fig_real_cases031.pdf}
\caption{The real-input case and post-hoc synthetic examples delimit the position
endpoint. Smaller centre errors and higher position F1 indicate better
localisation, while the 0.25 m annotation cutoff leaves the 0.5 m
matching rule unchanged. The real production replay is one reused room
with unverified physical clock synchronisation, 0.244 is the chair IoU,
and the synthetic examples come from a separate frozen 32-record batch.}
\end{figure*}

\section{One real-room feasibility
case}\label{one-real-room-feasibility-case}

The reused dormitory measures 2.84 by 4.045 m with a 2.09 m ceiling. Its
original RoNIN and phone-activity production chain retains the declared
room-specific relations and interaction protocol. Centre errors are
0.1117 m for bed, 0.1079 m for desk, 0.0907 m for cabinet and 0.2333 m
for chair. The inferred room registration is 170 degrees. The chair
footprint IoU is approximately 0.244, illustrating the distinction
between centre and footprint recovery.

A watch-driven cabinet replay retains four centre hits under seven
software clock shifts within plus or minus two seconds. Physical clock
synchronisation remains unverified. This n = 1 development-room replay
establishes feasibility for its production chain. The synthetic
displacement regressor is evaluated in its own generation domain.

\section{Discussion and conclusion}\label{discussion-and-conclusion}

The strongest controlled finding concerns contact placement. At the same
route and fourteen interactions, dispersed contacts increase whole-room
centre recovery by 24.92 percentage points over 10,000 generated
layouts. The activity map records the resulting partial and whole-room
disclosure under a specified protocol. Changing a declared prior, finite
decoder or scene family changes the operating point, which limits the
portability of a numerical activity budget.

The device results additionally depend on supervised synthetic dynamics,
posture semantics and event association. The two side-information
ablations establish that observations benefit the evaluated decoder,
while stronger prior-only baselines are needed to assess the full value
of side information. Fixed-record controls expose selection-dependent
refinement, but the replicated grid provides limited evidence for a
visit-by-error interaction. The primary new-versus-old whole-room
difference also remains statistically inconclusive at p = 0.1094. A
general raw-IMU activity threshold requires additional independently
controlled measurements.

Archived configurations, model metadata and commands document the
evaluated chains. Some historical caches retain only aggregate scores,
so independent box-level checks require replaying the saved inputs.
Across the controlled geometric experiments, device evaluation within
the training generator, and real-room case, acquisition, refinement and
finite selection remain practical sources of recovery loss.

\FloatBarrier
\section*{LLM Usage Statement}
OpenAI GPT models accessed through ChatGPT and Codex were used for
manuscript drafting and revision, literature and citation checking,
programming, figure preparation, and experiment design and result
auditing. Some code, plotting and auditing tasks were executed through
multi-step agent workflows, with outputs reviewed during the research
process. The authors are responsible for the submitted content.
{\small
\begin{thebibliography}{99}
\bibitem{r1} C. Mura, R. Pajarola, K. Schindler, and N. Mitra. Walk2Map: Extracting
floor plans from indoor walk trajectories. Computer Graphics Forum,
40:375--388, 2021.

\bibitem{r2} S. Basu, C. Amballa, Z. Xu, J. V. Sampedro, S. Nelakuditi, and R. R.
Choudhury. Contrastive diffusion guidance for spatial inverse problems.
arXiv:2509.26489, 2025.

\bibitem{r3} S. Ye, Y. Wang, J. Li, D. Park, C. K. Liu, H. Xu, and J. Wu. Scene
synthesis from human motion. In SIGGRAPH Asia Conference Papers, 2022.

\bibitem{r4} H.-Y. Hsu, T. Cheng, J. Wen, A. G. Schwing, and S. Wang. IMU4D: 4D
human-object understanding from wearable IMUs. arXiv:2604.21926, 2026.

\bibitem{r5} M. Hardegger, D. Roggen, S. Mazilu, and G. Troster. ActionSLAM: Using
location-related actions as landmarks in pedestrian SLAM. In IPIN, 2012.

\bibitem{r6} S. Herath, H. Yan, and Y. Furukawa. RoNIN: Robust neural inertial
navigation in the wild. In ICRA, pages 3146--3152, 2020.

\bibitem{r7} E. Kolve et al.~AI2-THOR: An interactive 3D environment for visual AI.
arXiv:1712.05474, 2017.

\bibitem{r8} M. Deitke et al.~ProcTHOR: Large-scale embodied AI using procedural
generation. arXiv:2206.06994, 2022.

\bibitem{r9} L. Breiman. Random forests. Machine Learning, 45:5--32, 2001.

\bibitem{r10} P. Geurts, D. Ernst, and L. Wehenkel. Extremely randomized trees.
Machine Learning, 63:3--42, 2006.

\bibitem{r11} E. B. Wilson. Probable inference, the law of succession, and statistical
inference. Journal of the American Statistical Association, 22:209--212,
1927.

\bibitem{r12} L. Lin, S. Xia, Z. Lai, L. Sun, J. Yang, and L. Pei. IMU-HOI: A
symbiotic framework for coherent human-object interaction and motion
capture via contact-conscious inertial fusion. In CVPR, pages
42901--42910, 2026.

\bibitem{r13} R. Yonetani and K. Hara. Map as a by-product: Collective landmark
mapping from IMU data and user-provided texts in situated tasks. PACM
IMWUT, 9(3):146, 2025.
\end{thebibliography}
}
\end{document}